\needspace{10\baselineskip}
\begin{center}
\pdfbookmark[1]{2. On Art. 53. Projection of a planetary orbit onto the ecliptic.}{zusatz-002}
[2. On Art. 53.]
\medskip

Projection of a planetary orbit onto the ecliptic.
\end{center}
\noindent \(e,\varphi,a,i,\AscendingNode,E\) with their usual meanings\\
\(\tilde\omega\) perihelion\\
\(g=\tilde\omega-\AscendingNode\)\\
\(\lambda\) hel[iocentric] l[ong]itude\\
\(\rho\) curtate distance [from the Sun].

[For the rectangular coordinates in the plane of the ecliptic, taking the line of nodes as the axis of abscissae, we have, by Arts. 53 and 58 of the Theoria motus corporum coelestium]
\[
\begin{aligned}
\text{[1]}\quad &a\cos g(\cos E-e)-a\cos\varphi\sin g\sin E\\
&\qquad=\rho\cos(\lambda-\AscendingNode)\\
&a\cos i\sin g(\cos E-e)+a\cos\varphi\cos i\cos g\sin E\\
&\qquad=\rho\sin(\lambda-\AscendingNode).
\end{aligned}
\]

[Put
\[
[2]\quad\begin{aligned}
a\cos g&=\alpha\cos\epsilon\cos\gamma-\beta\sin\epsilon\sin\gamma\\
a\cos\varphi\sin g&=\alpha\sin\epsilon\cos\gamma+\beta\cos\epsilon\sin\gamma\\
a\cos i\sin g&=\alpha\cos\epsilon\sin\gamma+\beta\sin\epsilon\cos\gamma\\
a\cos\varphi\cos i\cos g&=-\alpha\sin\epsilon\sin\gamma+\beta\cos\epsilon\cos\gamma,]
\end{aligned}
\]
\SourceLeaf{298}[thus] determine \(\alpha,\beta,\gamma,\epsilon\) by the equations
\[
[3]\quad\begin{aligned}
a(1-\cos\varphi\cos i)\cos g&=(\alpha-\beta)\cos(\gamma-\epsilon)\\
a(\cos i-\cos\varphi)\sin g&=(\alpha-\beta)\sin(\gamma-\epsilon)\\
a(1+\cos\varphi\cos i)\cos g&=(\alpha+\beta)\cos(\gamma+\epsilon)\\
a(\cos i+\cos\varphi)\sin g&=(\alpha+\beta)\sin(\gamma+\epsilon),
\end{aligned}
\]
[then we obtain
\[
\begin{aligned}
&\alpha(\cos(E+\epsilon)-e\cos\epsilon)\cos\gamma\\
&\qquad-\beta(\sin(E+\epsilon)-e\sin\epsilon)\sin\gamma
=\rho\cos(\lambda-\AscendingNode)\\
&\alpha(\cos(E+\epsilon)-e\cos\epsilon)\sin\gamma\\
&\qquad+\beta(\sin(E+\epsilon)-e\sin\epsilon)\cos\gamma
=\rho\sin(\lambda-\AscendingNode)
\end{aligned}
\]
or]
\[
[4]\quad\begin{aligned}
\alpha\cos(E+\epsilon)-\alpha e\cos\epsilon&=\rho\cos(\lambda-\gamma-\AscendingNode)\\
\beta\sin(E+\epsilon)-\beta e\sin\epsilon&=\rho\sin(\lambda-\gamma-\AscendingNode).
\end{aligned}
\]

Here \(\alpha,\beta\) are the two principal semiaxes of the projected ellipse;
\(-\alpha e\cos\epsilon,\allowbreak -\beta e\sin\epsilon\) the coordinates of its centre,
[\(\rho\cos(\lambda-\gamma-\AscendingNode),\allowbreak\rho\sin(\lambda-\gamma-\AscendingNode)\) the rectangular coordinates of the planet,] taking as the axis of abscissae the [parallel] to the axis \(2\alpha\) [drawn through the Sun]; \(\gamma+\AscendingNode\) its inclination to the equinoctial line.

The following auxiliary quantities may also be introduced in the calculation:
\[
\cos i=\operatorname{tang}k,\qquad
\cos\varphi=\operatorname{tang}l,\qquad
\frac{a}{\cos k\cos l}=c;
\]
then
\[
[5]\quad\begin{aligned}
a(1-\cos\varphi\cos i)&=c\cos(k+l)\\
a(\cos i-\cos\varphi)&=c\sin(k-l)\\
a(1+\cos\varphi\cos i)&=c\cos(k-l)\\
a(\cos i+\cos\varphi)&=c\sin(k+l).
\end{aligned}
\]

Example. Juno. Elements for 1810.
\[
\begin{aligned}
\varphi&=14^0\,45'\,35''\\
\log a&=0{,}42648\\
\text{trop. daily mot.}&=813''{,}5165\quad\text{sid.}=813''{,}379\\
i&=13^0\,4'\,12''\\
\AscendingNode&=171\quad7\quad48\\
\tilde\omega&=53\quad10\quad52
\end{aligned}
\]
\[
\begin{gathered}
\text{Epoch [of mean longitude] }1810=95^0\,33'\,50'',\\
g=242^0\,3'\,4''
\end{gathered}
\]
\SourceLeaf{299}
\[
\begin{array}{lcl}
\cos i&.\,.\,.\,.\,.\,.\,.\,.&9{,}988\,6011\\
\cos\varphi&.\,.\,.\,.\,.\,.\,.&9{,}985\,4277\\
\multicolumn{3}{l}{k=44^0\,14'\,53''{,}40}\\
\multicolumn{3}{l}{l=44\quad2\quad20{,}14}\\
\cos k&.\,.\,.\,.\,.\,.\,.&9{,}855\,1097\\
\cos l&.\,.\,.\,.\,.\,.\,.\,.&9{,}856\,6489\\
\hline
&&9{,}711\,7586\\
&&0{,}42648\\
\cline{3-3}
c&.\,.\,.\,.\,.\,.\,.\,.\,.\,.\,.&0{,}71472
\end{array}
\]
\[
\left[\begin{array}{l}k+l=\\k-l=\end{array}\right]
\begin{array}{l}88^0\,17'\,13''{,}54\\0\quad12\quad33{,}26\end{array}
\]
\[
\begin{array}{lcl}
\cos(k+l)&.\,.\,.\,.&8{,}47555\\
\sin(k+l)&.\,.\,.\,.&9{,}99981\\
\cos(k-l)&.\,.\,.\,.&0{,}00000\\
\sin(k-l)&.\,.\,.\,.&7{,}56252\\
\cos g&.\,.\,.\,.\,.\,.\,.\,.&9{,}67088_n\\
\sin g&.\,.\,.\,.\,.\,.\,.\,.&9{,}94614_n
\end{array}
\]
\[
\begin{array}{c}
\left[\begin{array}{l}
(\alpha-\beta)\cos(\gamma-\epsilon)\\
(\alpha-\beta)\sin(\gamma-\epsilon)\\
(\alpha+\beta)\cos(\gamma+\epsilon)\\
(\alpha+\beta)\sin(\gamma+\epsilon)
\end{array}\right]
\begin{array}{cl}
.\,.&8{,}86115_n\\.\,.&8{,}22338_n\\.\,.&0{,}38560_n\\.\,.&0{,}66067_n
\end{array}\\
\hline
\end{array}
\]
\[
\left[\begin{array}{l}\gamma-\epsilon=\\\gamma+\epsilon=\end{array}\right]
\begin{array}{l}192^0\,58'\,2''\\242\quad2\quad26\end{array}
\]
\[
\begin{array}{lcl}
\alpha-\beta&.\,.\,.\,.\,.\,.\,.\,.\,.\,.\,.&8{,}87237\\
\alpha+\beta&.\,.\,.\,.\,.\,.\,.\,.\,.\,.\,.&0{,}71457
\end{array}
\]
\[
\begin{aligned}
\gamma&=217^0\,30'\,14''\\
\epsilon&=24\quad32\quad12\\
\gamma+\AscendingNode&=28\quad38\quad2
\end{aligned}
\]
\[
\begin{array}{lcl}
\alpha&.\,.\,.\,.\,.\,.\,.\,.\,.\,.\,.\,.\,.\,.&0{,}41974\\
\beta&.\,.\,.\,.\,.\,.\,.\,.\,.\,.\,.\,.\,.\,.&0{,}40725\\
e&.\,.\,.\,.\,.\,.\,.\,.\,.\,.\,.\,.\,.\,.&9{,}40614\\
\hline
\end{array}
\]
\[
\begin{array}{c}
\left[\begin{array}{l}
\alpha e\\\cos\epsilon\\\beta e\\\sin\epsilon
\end{array}\right]
\begin{array}{cl}
.\,.\,.\,.\,.\,.\,.\,.\,.\,.&9{,}82588\\
.\,.\,.\,.\,.\,.\,.\,.\,.\,.&9{,}95890\\
.\,.\,.\,.\,.\,.\,.\,.\,.\,.&9{,}81339\\
.\,.\,.\,.\,.\,.\,.\,.\,.\,.&9{,}61834
\end{array}\\
\hline
\end{array}
\]
\[
\left.\begin{aligned}
-\alpha e\cos\epsilon&=-0{,}60923\\
-\beta e\sin\epsilon&=-0{,}27022
\end{aligned}\right\}
\begin{array}{l}\text{Coord.}\\\text{of c[entre].}\end{array}
\]

[The heliocentric coordinates can therefore now be calculated by the formula
\[
\begin{aligned}
&\rho\cos(\lambda-28^0\,38'\,2'')\\
&\qquad=-0{,}60923+(0{,}41974)\cos(E+24^0\,32'\,12'')\\
&\rho\sin(\lambda-28\quad38\quad2)\\
&\qquad=-0{,}27022+(0{,}40725)\sin(E+24\quad32\quad12)
\end{aligned}
\]
as soon as the eccentric anomaly is known.

The numbers in parentheses denote logarithms.]
\begin{center}\rule{.3\linewidth}{.4pt}\end{center}

The main steps in solving the inverse problem of finding a planetary orbit itself from the projections of the orbit and its focus are as follows:

\SourceLeaf{300}Let \(\alpha\) be the semimajor, \(\beta\) the semiminor axis of the projected ellipse; further, measured from the centre of the ellipse, relative to an arbitrary axis of abscissae
\begin{center}
\begin{tabular}{@{}lcl@{}}
\(P\)&.\ .\ .\ .&direction of \(\alpha\)\\
\(Q\)&.\ .\ .\ .&\begin{minipage}[t]{.75\linewidth}direction of the line to the focus, whose length \(=\delta\).\end{minipage}
\end{tabular}
\end{center}

Put
\[
[6]\quad\operatorname{tang}(R-P)\,.\,\operatorname{tang}(Q-P)
=-\frac{\beta\beta}{\alpha\alpha}\mathpunct{\raisebox{.45ex}{,}}
\]
[so that \(R\) is the direction of the diameter conjugate to the line to the focus; if \(q\) denotes the semidiameter passing through the focus and \(r\) its conjugate, then]
\[
[7]\quad\begin{aligned}
\BareRoot{\left\{\frac{\cos(Q-P)^2}{\alpha\alpha}+
\frac{\sin(Q-P)^2}{\beta\beta}\right\}}&=\frac{1}{q}\\
\BareRoot{\left\{\frac{\cos(R-P)^2}{\alpha\alpha}+
\frac{\sin(R-P)^2}{\beta\beta}\right\}}&=\frac{1}{r}\mathpunct{\raisebox{.45ex}{,}}
\end{aligned}
\]
therefore
\[
[8]\quad\frac{qq}{rr}=-\frac{\sin2(R-P)}{\sin2(Q-P)}\mathpunct{\cdot}
\]

[The eccentricity of the ellipse is]
\[
[9]\quad e=\frac{\delta}{q}\mathpunct{\cdot}
\]

[By introducing the angle \(\epsilon\), also used above, which represents the eccentric anomaly corresponding to the radius vector \(q\) in the projected ellipse, the endpoints of the conjugate radius vectors \(q\) and \(r\) can be expressed in rectangular coordinates, now referred to the axis \(2\alpha\) as the axis of abscissae, as follows:
\[
\begin{aligned}
q\cos(Q-P)&=\alpha\cos\epsilon,&
r\cos(R-P)&=-\alpha\sin\epsilon,\\
q\sin(Q-P)&=\beta\sin\epsilon,&
r\sin(R-P)&=\beta\cos\epsilon;
\end{aligned}
\]
if from these we derive the rectangular coordinates referred to the line of nodes as the axis of abscissae, by rotating through the angle \(\gamma=P-\AscendingNode\), then, in view of equations 2, we have:]
\[
[10]\quad\begin{aligned}
q\cos(Q-\AscendingNode)&=a\cos g,\\
\frac{r}{\BareRoot{(1-ee)}}\cos(R-\AscendingNode)&=-a\sin g\\
q\sin(Q-\AscendingNode)&=a\cos i\sin g,\\
\frac{r}{\BareRoot{(1-ee)}}\sin(R-\AscendingNode)&=a\cos i\cos g,
\end{aligned}
\]
\SourceLeaf{301}[from which we readily derive]
\[
[11]\quad\frac{qq}{rr}(1-ee)
=-\frac{\sin2(R-\AscendingNode)}{\sin2(Q-\AscendingNode)}\mathpunct{\cdot}
\]

[Further we find]
\[
[12]\quad\begin{aligned}
qq+\frac{rr}{1-ee}&=aa(1+\cos i^2)\\
\frac{qr}{\BareRoot{(1-ee)}}\sin(R-Q)&=aa\cos i\\
\frac{qr}{\BareRoot{(1-ee)}}\cos(R-Q)&=-\tfrac{1}{2}aa\sin i^2\sin2g\\
qq-\frac{rr}{1-ee}&=aa\sin i^2\cos2g.
\end{aligned}
\]

[In a similar way we derive
\[
\begin{aligned}
\frac{qr}{\BareRoot{(1-ee)}}\cos(Q+R-2\AscendingNode)
&=-\tfrac{1}{2}aa(1+\cos i^2)\sin2g\\
\frac{qr}{\BareRoot{(1-ee)}}\sin(Q+R-2\AscendingNode)
&=aa\cos i\cos2g,
\end{aligned}
\]
therefore, in view of equations 12:]
\[
[13]\quad\begin{aligned}
&aa\sin i^2\cos(Q+R-2\AscendingNode)\\
&\qquad=\left(qq+\frac{rr}{1-ee}\right)\cos(R-Q)\\
&aa\sin i^2\sin(Q+R-2\AscendingNode)\\
&\qquad=\left(qq-\frac{rr}{1-ee}\right)\sin(R-Q).
\end{aligned}
\]

[The equations derived contain everything needed to solve the problem and can be used in various ways; if \(\alpha,\beta,P,Q,\delta\) are given, one can, for example, determine \(R\) from 6, \(q\) and \(r\) from 7, \(e\) from 9; then \(\AscendingNode,a,i,g\) are found from 13 and 10. The position of the axis of abscissae from which the directions \(P,Q,R,\AscendingNode\) are measured remains entirely arbitrary.]
\begin{center}\rule{.3\linewidth}{.4pt}\end{center}

